% Historical source (removed by author request): references/_OLD_Übungsaufgaben/Übungsblatt 1 - Grundlagen/Blatt1.tex:34
\begin{ExerciseBlock}
\ExerciseHeading{02}{Mengenalgebra}{fundamentals_set_operations}
Gegeben sind die Mengen.
\[
\begin{aligned}
	A &= \left\lbrace 1, 2, 3 \right\rbrace \\
	B &= \left\lbrace 3, 4, 5, 6, 7, 8, 9, 10 \right\rbrace \\
	C &= \left\lbrace x \in \N \mid x\text{ ist eine Primzahl} \right\rbrace \\
	D &= \left\lbrace 2k \mid k \in \N \right\rbrace
\end{aligned}
\]

Berechnen Sie:
\begin{enumerate}[label=\alph*)]
	\item $\displaystyle  A \cap C = $
	\item $\displaystyle A \cup B = $
	\item $\displaystyle \left( A \cup B \right) \setminus C = $
	\item $\displaystyle D \cap C = $
	\item $\displaystyle D \cap B = $
	\item $\displaystyle C \cap C = $
	\item $\displaystyle C \cup C = $
	\item $\displaystyle C \setminus C = $
	\item $\displaystyle D \cap \emptyset = $
	\item $\displaystyle D \cup \emptyset = $
\end{enumerate}

\end{ExerciseBlock}
\UseExerciseSolution{fundamentals_set_operations}
